% This is part of Mes notes de mathématique
% Copyright (c) 2008-2019
%   Laurent Claessens, Carlotta Donadello
% See the file fdl-1.3.txt for copying conditions.


%+++++++++++++++++++++++++++++++++++++++++++++++++++++++++++++++++++++++++++++++++++++++++++++++++++++++++++++++++++++++++++
\section{Formes bilinéaires et quadratiques}
%+++++++++++++++++++++++++++++++++++++++++++++++++++++++++++++++++++++++++++++++++++++++++++++++++++++++++++++++++++++++++++

Plus à propos de formes bilinéaires dans le thème \ref{THEMEooOAJKooEvcCVn}.

%---------------------------------------------------------------------------------------------------------------------------
\subsection{Généralités}
%---------------------------------------------------------------------------------------------------------------------------

\begin{definition}[\cite{RUAoonJAym}]   \label{DefBSIoouvuKR}
    Soit un espace vectoriel \( E\) et \( \eF\) un corps de caractéristique différente de \( 2\). Une \defe{forme quadratique}{forme!quadratique} sur \( E\) est une application \( q\colon E\to \eF\) pour laquelle il existe une forme bilinéaire symétrique \( b\colon E\times E\to \eF\) satisfaisant \( q(x)=b(x,x)\) pour tout \( x\in E\).

    L'ensemble des formes quadratiques réelles sur \( E\) est noté \( Q(E)\)\nomenclature[B]{\( Q(E)\)}{formes quadratiques réelles sur \( E\)}.
\end{definition}

\begin{proposition} \label{PROPooZLXVooOsXCcB}
    Soit une forme bilinéaire \( b\) et la forme quadratique associée \( q\). Alors nous avons l'\defe{identité de polarisation}{identité de polarisation} :
\begin{equation}    \label{EqMrbsop}
    b(x,y)=\frac{ 1 }{2}\big( q(x)+q(y)-q(x-y) \big).
\end{equation}
\end{proposition}

\begin{proof}
    Il suffit de substituer dans le membre de droite \( q(x)=b(x,y)\) et d'utiliser la bilinéarité :
    \begin{subequations}
        \begin{align}
            q(x)+q(y)-q(x-y)&=b(x,x)+b(y,y)-b(x-y,x-y)\\
            &=b(x,x)+b(y,y)-b(x)+b(x,y)+b(y,x)-b(y,y)\\
            &=2b(x,y)
        \end{align}
    \end{subequations}
    où nous avons utilisé le fait que \( b\) est symétrique : \( b(x,y)=b(y,x)\).
- \end{proof}
- 

\begin{lemma}       \label{LEMooLKNTooSfLSHt}
    Si \( q\) est une forme quadratique, il existe une unique forme bilinéaire \( b\) telle que \( q(x)=b(x,x)\).
\end{lemma}

\begin{proof}
    L'existence n'est pas en cause : c'est la définition d'une forme quadratique. Pour l'unicité, étant donné une forme quadratique, la forme bilinéaire \( b\) doit forcément vérifier l'identité de polarisation de la proposition \ref{PROPooZLXVooOsXCcB}. Elle est donc déterminée par \( q\).
\end{proof}
Notons la division par \( 2\) qui est le pourquoi de la demande de la caractéristique différente de \( 2\) pour \( \eF\) dans la définition de forme quadratique.

%---------------------------------------------------------------------------------------------------------------------------
\subsection{Matrice associée à une forme bilinéaire}
%---------------------------------------------------------------------------------------------------------------------------
 
\begin{definition}      \label{DEFooAOGPooXWXUcN}
    Soit une forme bilinéaire \( b\colon E\times E\to \eK\) et une base quelconque \( \{ f_{\alpha} \}\) de \( E\). Nous définissons les nombres
    \begin{equation}    \label{EQooCUGFooRlKUtu}
        B_{\alpha\beta}=b(f_{\alpha},f_{\beta}),
    \end{equation}
    qui forment une matrice symétrique dans \( \eM(n,\eK)\). Cette matrice est la \defe{matrice associée}{matrice d'une forme bilinéaire} à la forme bilinéaire \( b\).

    La matrice d'une forme quadratique est celle associée à sa forme bilinéaire associée.
\end{definition}

Alors nous avons aussi, si \( x=\sum_{\alpha}x_{\alpha}f_{\alpha}\) et \( y=\sum_{\beta}y_{\beta}f_{\beta}\) :
  \begin{equation}        \label{EQooQFMWooVKVLMx}
    b(x,y)=\sum_{\alpha\beta}x_{\alpha}y_{\beta}b(f_{\alpha},f_{\beta})=\sum_{\alpha\beta}B_{\alpha\beta}x_{\alpha}y_{\beta}.
  \end{equation}
  La matrice \( B\) est la matrice de \( b\) dans la base \( \{ f_{\alpha} \}\) de \( V\).
  
Nous pouvons aussi écrire cela avec un produit scalaire. En utilisant la convention \eqref{EQooAXRJooUwHbjB} et les choses autour (voir aussi \ref{SECooBTTTooZZABWA}),
\begin{equation}
    b(x,y)=\sum_{\alpha}x_{\alpha}\sum_{\beta}B_{\alpha\beta}y_{\beta}=\sum_{\alpha}x_{\alpha}(By)_{\alpha}=x\cdot By.
\end{equation}
où le point est le produit scalaire usuel (composante par composante).

\begin{normaltext}
    De nombreux auteurs préfèrent écrire des choses comme \( x^tBy\) ou \( xB^ty\) ou \( xBy^t\) et se poser de longues questions sur qui est un «vecteur colonne» et qui est un «vecteur ligne», et si la matrice \( B\) soit être transposée ou non. Toutes ces notations servent(?) à cacher un bête produit scalaire.
\end{normaltext}

\begin{normaltext}
    Notons que la matrice associée à une forme bilinéaire (ou quadratique associée) est uniquement valable pour une base donnée. Si nous changeons de base, la matrice change. Cependant lorsque nous travaillons sur \( \eR^n\), la base canonique est tellement canonique que nous allons nous permettre de parler de «la» matrice associée à une forme bilinéaire.
\end{normaltext}

%---------------------------------------------------------------------------------------------------------------------------
\subsection{Changement de base : matrice d'une forme bilinéaire}
%---------------------------------------------------------------------------------------------------------------------------

\begin{proposition}[Voir la section \ref{SECooBTTTooZZABWA}]     \label{PROPooLBIOooUpzxXA}
    
    Soit une forme bilinéaire\footnote{Définition~\ref{DEFooEEQGooNiPjHz}} \( b\colon V\times V\to \eK\) dont la matrice\footnote{Définition~\ref{EQooCUGFooRlKUtu}.} dans la base \( \{ e_i \}\) est \( A\) et celle dans la base \( \{ f_{\alpha} \}\) est \( B\). Nous supposons que les bases sont liées par \( f_{\alpha}=\sum_{i}Q_{i\alpha}e_i\). Alors
\begin{equation}        \label{EQooZUVTooKjqnJj}
    B=Q^tAQ.
\end{equation}
\end{proposition}

\begin{proof}
    Soit \( x,x'\in V\) de coordonnées \( (x_i)\) et \( (x'_i)\) dans la base \( \{ e_i \}\) et \( (y_{\alpha})\), \( (y'_{\alpha})\) dans la base \( \{ f_{\alpha} \}\). Par définition de la matrice associée à une forme bilinéaire,
    \begin{equation}
        b(x,x')=\sum_{ij}A_{ij}x_ix'_j=\sum_{\alpha\beta}B_{\alpha\beta}y_{\alpha}y'_{\beta}.
    \end{equation}
    En remplaçant les \( x_i\) et \( x'_i\) par leurs valeurs en fonction de \( y_{\alpha}\) et \( y'_{\beta}\) données par la proposition \ref{PROPooNYYOooHqHryX}, nous trouvons
    \begin{subequations}
        \begin{align}
            b(x,x')&=\sum_{ij\alpha\beta}A_{ij}Q_{i\alpha}y_{\alpha}Q_{j\beta}y'_{\beta}\\
            &=\sum_{\alpha\beta}(Q^tAQ)_{\alpha\beta}y_{\alpha}y'_{\beta}
        \end{align}
    \end{subequations}
    où \( Q^t\) désigne la transposée de la matrice \( Q\) :  \( Q^t_{ij}=Q_{ji}\). Vu que les nombres \( y_{\alpha}\) et \( y'_{\beta}\) sont arbitraires nous déduisons\footnote{Lemme~\ref{LEMooLXAHooPRyHaF}.} que \( B=Q^tAQ\).
\end{proof}

\begin{remark}      \label{REMooNEJLooSqgeih}
    Notons que cette «loi de transformation» n'est pas la même que celle pour une application linéaire\footnote{Proposition \ref{PROPooNZBEooWyCXTw}.}. Ici nous avons \( Q^t\) alors que pour les applications linéaires nous avions \( Q^{-1}\).

    Pour cette raison, tant que nous travaillons avec des bases orthonormées, c'est-à-dire tant que \( Q\) est orthogonale\footnote{Définition~\ref{DefMatriceOrthogonale}.}, nous pouvons confondre une application linéaire avec une application bilinéaire en passant par la matrice. Mais cette identification n'est pas du tout canonique : elle repose sur le fait que les bases soient orthonormées.

    Il en découle que la réduction des endomorphismes et la réduction des formes bilinéaires ne sont pas tout à fait les mêmes théories. Par exemple la pseudo-diagonalisation simultanée (corolaire~\ref{CorNHKnLVA}) est un résultat de réduction de forme bilinéaire et non d'endomorphismes.
\end{remark}

\begin{proposition}     \label{PROPooCIEUooODqfwm}
    Soit une forme quadratique \( q\colon E\to \eK\) et sa matrice\footnote{Matrice associée à une forme quadratique, définition \ref{DEFooAOGPooXWXUcN}.} \( (q_{ij})\in \eM(n,\eK)\). Nous avons
    \begin{equation}        \label{EQooQROVooBfshbZ}
        q(x)=\sum_{i=1}^n\sum_{j=1}^nq_{ij}x_ix_j=\sum_{i=1}^nq_{ii}x_i^2+2\sum_{1\leq i <j\leq n}q_{ij}x_ix_j.
    \end{equation}
\end{proposition}

%--------------------------------------------------------------------------------------------------------------------------- 
\subsection{Réduction de Gauss}
%---------------------------------------------------------------------------------------------------------------------------

\begin{theorem}[Réduction de Gauss\cite{BIBooNUUEooJUjLpy,BIBooUULNooUtlrar}]     \label{THOooOMMFooKxqICS}
    Soit une forme quadratique non nulle \( q\) sur l'espace vectoriel \( E\) sur le corps \( \eK\). Il existe des formes linéaires linéairement indépendantes \(\{ l_i \}_{i=1,\ldots, n}\) et des coefficients \( \alpha_i\in \eK\) tels que 
        \begin{equation}
            q(x)=\sum_{i=1}^n\alpha_il_i(x)^2.
        \end{equation}
\end{theorem}

\begin{proof}
    Notre point de départ est la formule \eqref{EQooQROVooBfshbZ} pour la forme quadratique. Nous allons faire la preuve par récurrence sur la dimension de l'espace. Si \( n=1\), alors nous avons seulement
    \begin{equation}
        q(x)=\alpha x^2
    \end{equation}
    et donc le théorème est fait avec \( l(x)=x\).

    Nous supposons que le théorème est prouvé pour tout espace de dimension \( n\). Une forme quadratique pour un espace de dimension \( n+1\) s'écrit
    \begin{equation}
        q(x)=\sum_{i=1}^{n+1}q_{ii}x_i^2+2\sum_{1\leq i < j\leq n+1}q_{ij}x_ix_j.
    \end{equation}
    Vu que \( q\) est non nulle, un des \( q_{ij}\) est non nul. Nous allons diviser en plusieurs cas.
    \begin{itemize}
        \item
            \( q_{11}\neq 0\)
        \item
            \( q_{kk}\neq 0\) avec \( k\neq 1\)
        \item
            \( q_{12}\neq 0\)
        \item
            \( q_{kl}\neq 0\) avec \( (k,l)\neq (1,2)\).
    \end{itemize}
    Ces cas ne sont pas exclusifs, mais ils couvrent toutes les possibilités.
\end{proof}


%--------------------------------------------------------------------------------------------------------------------------- 
\subsection{Orthogonalité}
%---------------------------------------------------------------------------------------------------------------------------

\begin{proposition}[\cite{BIBooUULNooUtlrar}]
    Soit un espace vectoriel \( (E,\eK)\) et une forme quadratique \( q\). Une base de \( E\) est \( q\)-orthogonale si et seulement si la matrice de \( q\) dans cette base est diagonale.
\end{proposition}


\begin{theorem}[\cite{BIBooUULNooUtlrar}]       \label{THOooIDMPooIMwkqB}
    Toute forme quadratique sur un espace vectoriel de dimension finie admet une base formée de vecteurs \( 2\) à \( 2\) orthogonaux (pour la forme considérée).
\end{theorem}
  
  %--------------------------------------------------------------------------------------------------------------------------- 
  \subsection{Diagonalisation}
  %---------------------------------------------------------------------------------------------------------------------------
  
Le théorème \ref{THOooIDMPooIMwkqB} a déjà donné une base orthogonale pour toute forme quadratique sur un espace vectoriel \( (E,\eK)\) de dimension finie. Dans le cas de \( \eR^n\), nous pouvons en donner une preuve basée sur le théorème spectral, c'est la proposition \ref{PROPooUKRUooGRIDHt}.

\begin{proposition}     \label{PROPooUKRUooGRIDHt}
    Soit une forme bilinéaire symétrique \( b\) sur un \( \eR^n\). Il existe une matrice orthogonale \( Q\) telle que 
      \begin{enumerate}
          \item
              \( D=Q^tbQ\) est diagonale
          \item
              \( D(x,y)=b(Qx,Qy)\) pour tout \( x,y\in E\).
      \end{enumerate}
    Il existe une base \( (f_i)_{i=1,\ldots, n}\) qui est \( b\)-orthogonale.
\end{proposition}

      Dans cet énoncé, nous mélangeons sans vergogne les formes et les matrices, en supposant qu'une base soit fixée\footnote{Autrement dit, si vous avez en tête d'utiliser cette proposition pour \( \eR^n\) c'est bon; mais sinon vous devez choisir une base et considérer toutes les matrices dans cette base.}. Par exemple
    \begin{equation}
        D(x,y)=\sum_{ij}D_{ij}x_iy_j.
    \end{equation}


\begin{proof}
    Pour la matrice diagonale, c'est le théorème spectral \ref{ThoeTMXla}\ref{ITEMooMWWRooXxGONW} qui joue parce que la matrice d'une forme bilinéaire symétrique est symétrique (c'est vu de la définition \eqref{EQooCUGFooRlKUtu}).

    Pour le reste c'est un calcul :
    \begin{subequations}
        \begin{align}
            D(x,y)&=\sum_{ijkl}Q^t_{ik}b_{kl}Q_{lj}x_iy_j\\
            &=\sum_{ijkl}b_{kl}(Q_{ki}x_i)(Q_{lj}y_j)\\
            &=\sum_{kl}b_{kl}(Qx)_k(Qy)_l\\
            &=b(Qx,Qy).
        \end{align}
    \end{subequations}
    Nous avons utilisé le produit matrice fois vecteur donné par \eqref{EQooQFVTooMFfzol}.
    
         En ce qui concerne l'existence d'une base \( b\)-orthogonale, vu que \( D\) est diagonale, nous avons, pour \( i\neq j\) que \( D(e_i,e_j)=0\). Donc en posant \( f_i=Qe_i\), nous trouvons
     \begin{equation}
         0=D(e_i,e_j)=b(Qe_i,Qe_j)=b(f_i,f_j).
     \end{equation}
     La base \( (Qe_i)_{i=1,\ldots, n}\) est donc \( b\)-orthogonale.
\end{proof}

%--------------------------------------------------------------------------------------------------------------------------- 
\subsection{Isométrie, forme quadratique et bilinéaire}
%---------------------------------------------------------------------------------------------------------------------------

\begin{example}
    La forme quadratique \( q(x)=x_1^2+x_2^2\) donne la norme euclidienne. La forme bilinéaire associée est \( b(x,y)=x_1y_1+x_2y_2\), qui est le produit scalaire usuel.
\end{example}

Il ne faudrait pas déduire trop vite que la formule \( \| x \|^2=q(x)\) donne une norme dès que \( q\) est non dégénérée. En effet \( q\) peut ne pas être définie positive. La forme \( q(x)=x_1^2-x_2^2\) prend des valeurs positives et négatives. A fortiori \( d(x,y)=q(x-y)\) ne donne pas toujours une distance.

\begin{definition}      \label{DEFooECTUooRxBhHf}
    Une \defe{isométrie}{isométrie!de forme quadratique} pour la forme quadratique \( q\) est une application bijective \( f\colon V\to V\) telle que 
    \begin{equation}
     q(x-y)=q\big( f(x)-f(y) \big).
    \end{equation}
     Dans les cas où \( q\) donne une distance, alors c'est une isométrie au sens usuel.
\end{definition}

\begin{definition}[Thème \ref{THMooVUCLooCrdbxm}]      \label{DEFooIQURooMeQuqX}
    Soit un espace vectoriel \( E\) muni d'une forme bilinéaire \( b\). Une \defe{isométrie}{isométrie (forme bilinéaire)} pour \( b\) est une bijection \( f\colon E\to E\) telle que
    \begin{equation}
        b\big( f(x),f(y) \big)=b(x,y)
    \end{equation}
    pour tout \( x,y\in E\).
\end{definition}

\begin{lemma}   \label{LemewGJmM}
    Soient \( q\) une forme quadratique et \( b\) la forme bilinéaire associée par le lemme~\ref{LEMooLKNTooSfLSHt}. Une application \( f\colon E\to E\) telle que \( f(0)=0\) est une isométrie pour \( b\) si et seulement si elle est une isométrie pour \( q\).
\end{lemma}

\begin{proof}
    Pour une application bijective \( f\colon E\to E\) telle que \( f(0)=0\), nous devons prouver l'équivalence des propriétés suivantes :
    \begin{enumerate}
        \item
            \( b\big( f(x),f(y) \big)=b(x,y)\) pour tout \( x,y\in E\);
        \item
            \( q\big( f(x)-f(y) \big)=q(x-y)\) pour tout \( x,y\in E\).
    \end{enumerate}

    Dans le sens direct, en posant \( x=y\) nous trouvons tout de suite \( q(f(x))=q(x)\); ensuite en utilisant la distributivité de \( b\),
    \begin{subequations}
        \begin{align}
            q\big( f(x)-f(y) \big)&=b\big( f(x)-f(y),f(x)-f(y) \big)\\
            &=q\big( f(x) \big)-2b\big( f(x),f(y) \big)+q\big( f(y) \big)\\
            &=q(x)+q(y)-2b(x,y)\\
            &=q(x-y).
        \end{align}
    \end{subequations}

    Dans l'autre sens, nous commençons par remarquer que l'hypothèse \( f(0)=0\) donne \( q(x)=q\big( f(x) \big)\). Ensuite nous utilisons l'identité de polarisation \eqref{EqMrbsop} :
    \begin{subequations}
        \begin{align}
            b\big( f(x),f(y) \big)&=\frac{ 1 }{2}\big[ q\big( f(x) \big)+q\big( f(y) \big)-q\big( f(x-y) \big) \big]\\
            &=\frac{ 1 }{2}\big[ q(x)+q(y)-q(x-y) \big]\\
            &=b(x,y).
        \end{align}
    \end{subequations}
\end{proof}


%+++++++++++++++++++++++++++++++++++++++++++++++++++++++++++++++++++++++++++++++++++++++++++++++++++++++++++++++++++++++++++
\section{Conventions sur les matrices et changement de bases}
%+++++++++++++++++++++++++++++++++++++++++++++++++++++++++++++++++++++++++++++++++++++++++++++++++++++++++++++++++++++++++++
\label{SECooBTTTooZZABWA}

Nous nous proposons à présent de fixer toutes les notations concernant le calcul matriciel et les changements de bases\footnote{Concernant la matrice d'une différentielle, c'est la proposition \ref{PROPooBMROooThgLuU}. Nous n'en parlerons pas ici parce que c'est pour plus tard.}. Nous considérons des espaces vectoriels \( V\) et \( W\) respectivement munis de bases \( \{ e_i \}\) et \( \{ f_{\alpha} \}\). Ils sont tous deux sur le corps  commutatif \( \eK\). Si \( T\colon V\to W \) est une application linéaire, alors sa matrice est définie en la définition \ref{DEFooJVOAooUgGKme} (et les résultats qui montrent que le \( \psi\) est une bijection comme il faut). Ici nous prenons une notation plus détendue en notant \( T\) l'apllication linéaire et sa matrice. 

En premier lieu nous avons
\begin{equation}        \label{EQooOMSCooGsSBIA}
    T_{\alpha i}=T(e_i)_{\alpha}.
\end{equation}
Donc aussi
\begin{equation}
    T(e_i)=\sum_{\alpha}T_{\alpha i}f_{\alpha}.
\end{equation}
À ce point, nous avons l'impression de prendre une convention qui donne la matrice à l'envers par rapport à ce que l'on voudrait pour que les indices identiques se suivent. La raison est que la formule \eqref{EQooUELWooErOncF} de l'action d'une application linéaire sur un vecteur sera, elle, à l'endroit.

En ce qui concerne l'application d'une matrice à un vecteur, c'est définit en \eqref{EQooQFVTooMFfzol}; nous avons alors, pour \( x=\sum_ix_ie_i\) :
\begin{equation}
    T(x)=\sum_iT(e_i)=\sum_{i\alpha}T_{\alpha i}x_if_{\alpha}
\end{equation}
ou encore, en pour les composantes :
\begin{equation}        \label{EQooUELWooErOncF}
    T(x)_{\alpha}=\sum_{i}T_{\alpha i}x_i.
\end{equation}

Lorsque nous avons une base orthonormée nous avons aussi
\begin{equation}        \label{EQooDSKBooQkgtWv}
    A_{ij}=\langle Ae_j, e_i\rangle .
\end{equation}

Ici \( t\) est une application \( V\to V\) et \( A\in \eM(n,\eK)\) est simplement un tableau de nombres sans prétentions d'être une application linéaire. Cependant, \( A\) peut être vu comme application linéaire \( \eK\to \eK\), et l'ensemble de nombres \( \{ x_i \}_{i=1,\ldots, n}\) peut être vu comme vecteur de \( \eK^n\). Dans ce contexte nous pouvons écrire
\begin{equation}
    t(x)_i=(Ax)_i.
\end{equation}
Cela signifie que si on identifie un vecteur au vecteur de ses composantes, l'application linéaire se réduit au produit «matrice fois vecteur» sur le corps de base.

Tant que nous y sommes, nous avons aussi la formule suivante dans \( \eR^n\) muni du produit scalaire usuel :
\begin{equation}        \label{EQooFYNYooFQhVyE}
    x\cdot Ay=\sum_kx_k(Ay)_k=\sum_{kl}x_kA_{kl}y_l=\sum_{kl}A_{kl}x_ky_l.
\end{equation}
En voila au moins une pour laquelle les indices tombent au bons endroits.

%---------------------------------------------------------------------------------------------------------------------------
\subsection{Le changement de base}
%---------------------------------------------------------------------------------------------------------------------------

Soit un espace vectoriel \( V\) muni de deux bases \( \{ e_i\}_{i=1,\ldots, n}\) et \( \{ f_{\alpha}\}_{\alpha=1,\ldots, n}\). Les deux bases sont liées entre elles par
\begin{equation}        \label{EQooFRQRooSMsQQB}
    f_{\alpha}=\sum_iQ_{i\alpha}e_i.
\end{equation}
Ici \( Q\) n'est pas une application linéaire \( V\to V\) : \( Q\) est seulement un tableau de nombres, donnant les coordonnées des vecteurs \( f_{\alpha}\) dans la base de \( e_i\). Éventuellement \( Q\) peut être vu comme une application linéaire \( \eK^n\to \eK^n\).

Dans la suite nous nommerons \( Q^{-1}\) la matrice inverse de \( Q\). Inverse au sens des bêtes tableaux de nombres, sans interprétations en tant qu'application linéaire. De même pour \( Q^t\) qui est la transposée de \( Q\).

\begin{lemma}       \label{LEMooLXAHooPRyHaF}
    Soient des matrices \( A,B\in \eM(n,\eK)\). Si pour tout \( x,y\in \eK^n\) nous avons
    \begin{equation}
        \sum_{ij}A_{ij}x_iy_j=\sum_{ij}B_{ij}x_iy_j
    \end{equation}
    alors \( A=B\).
\end{lemma}

\begin{proof}
    Il suffit de choisir \( x_i=\delta_{ik}\) et \( y_j=\delta_{jl}\), et d'effectuer les sommes; par exemple
    \begin{equation}
        \sum_{ij}A_{ij}\delta_{ik}y_j=\sum_jA_{kj}y_j.
    \end{equation}
    Après avoir effectué toutes les sommes nous nous retrouvons avec \( A_{kl}=B_{kl}\), ce qui signifie \( A=B\).
\end{proof}

%---------------------------------------------------------------------------------------------------------------------------
\subsection{Changement de base : vecteurs de base}
%---------------------------------------------------------------------------------------------------------------------------

Nous multiplions l'égalité \eqref{EQooFRQRooSMsQQB} par \( Q^{-1}_{\alpha j}\) \footnote{Attention à la bonne interprétation de ce nombre : on fait bien référence à l'élément situé en \( (\alpha, j) \) de la matrice \( Q^{-1} \), et pas autre chose.} et nous sommons sur \( \alpha\) :
\begin{equation}
    \sum_{\alpha}Q^{-1}_{\alpha j}f_{\alpha}=\sum_{i\alpha}(A_{i\alpha}Q^{-1}_{\alpha j})e_i=e_j.
\end{equation}
Donc :
\begin{equation}    \label{EQooZQPAooAbKAdg}
    e_i=\sum_{\alpha}Q^{-1}_{\alpha i}f_{\alpha}.
\end{equation}

%---------------------------------------------------------------------------------------------------------------------------
\subsection{Changement de base : coordonnées}
%---------------------------------------------------------------------------------------------------------------------------

Soit un vecteur \( x\in V\). Il peut être écrit dans les deux bases :
\begin{equation}
    x=\sum_ix_ie_i=\sum_{\alpha}y_{\alpha}f_{\alpha}.
\end{equation}
En remplaçant \( e_i\) par sa valeur \eqref{EQooZQPAooAbKAdg} nous avons l'égalité
\begin{equation}
    \sum_{i\alpha}x_iQ^{-1}_{\alpha i}f_{\alpha}=\sum_{\alpha}y_{\alpha}f_{\alpha}.
\end{equation}
Vu que les \( f_{\alpha}\) sont linéairement indépendants, l'égalité des sommes donne l'égalité de chacun de termes :
\begin{equation}        \label{EQooFXYLooCRmRdA}
    y_{\alpha}=\sum_ix_iQ^{-1}_{\alpha i}.
\end{equation}
En identifiant \( x\in V\) au vecteur dans \( \eK^n\) de ses coordonnées dans la base \( \{ e_i \}\) nous pouvons écrire
\begin{equation}
    y_{\alpha}=(Q^{-1}x)_{\alpha},
\end{equation}
ou pire :
\begin{equation}
    y=Q^{-1}x.
\end{equation}
Cette dernière égalité repose sur un petit paquet d'abus de notations qu'il convient de bien comprendre. Ici, \( x\) et \( y\) sont les éléments de \( \eK^n\) donnés par les composantes de \( x\) dans les bases \( \{ e_i \}\) et \( \{ f_{\alpha} \}\), et \( Q\) est vu comme une matrice, un opérateur linéaire sur \( \eK^n\).

Une chose agréable avec cette façon d'écrire est que nous trouvons tout de suite la transformation inverse \( x=Qy\) qui peut être écrire de différentes manières :
\begin{subequations}
    \begin{align}
        x&=Qy\\
        x_i&=(Qy)_i   \label{SUBEQooPVGBooDafCcBk}  \\
        x_i&=\sum_{\alpha}Q_{i\alpha}y_{\alpha}
    \end{align}
\end{subequations}
Attention à l'ordre des indices dans la dernière égalité : la matrice \( Q\) vient avec les indices dans l'ordre \( i\alpha\), tandis que la matrice \( Q^{-1}\) vient avec les indices dans l'ordre opposé : \( \alpha i\). C'est pour cela qu'il est intéressant de noter avec des lettres latines les indices se rapportant à la première base et avec des lettres grecques ceux se rapportant à la seconde base.

%---------------------------------------------------------------------------------------------------------------------------
\subsection{Changement de base : matrice d'une application linéaire}
%---------------------------------------------------------------------------------------------------------------------------

\begin{proposition}     \label{PROPooNZBEooWyCXTw}
    Soit une application linéaire \( t\colon V\to V\) de matrices \( A\) et \( B\) dans les bases \( \{ e_i \}\) et \( \{ f_{\alpha} \}\). Si les bases sont liées par
    \begin{equation}
        f_{\alpha}=\sum_iQ_{i\alpha}e_i,
    \end{equation}
    alors les matrices \( A\) et \( B\) sont liées par
    \begin{equation}
        B=Q^{-1}AQ.
    \end{equation}
\end{proposition}

\begin{proof}
    L'hypothèse sur le fait que \( A\) et \( B\) sont les matrices de \( t\) signifie que pour tout \( x\in V\),
    \begin{equation}
        t(x)=\sum_{ij}A_{ji}x_ie_j=\sum_{\alpha\beta}B_{\alpha\beta}y_{\beta}f_{\alpha}.
    \end{equation}
    En remplaçant \( e_j\) par son expression \eqref{EQooZQPAooAbKAdg} en termes des \( f_{\alpha}\) et \( x_i\) par son expression \eqref{SUBEQooPVGBooDafCcBk}, nous avons
    \begin{subequations}
        \begin{align}
            (By)_{\alpha}&=\sum_{ij\alpha}A_{ji}(Qy)_iQ^{-1}_{\alpha j}f_{\alpha}\\
            &=\sum_{i \alpha}(Q^{-1}A)_{\alpha i}(Qy)_if_{\alpha}\\
            &=\sum_{\alpha}(Q^{-1} AQy)_{\alpha}f_{\alpha}.
        \end{align}
    \end{subequations}
    Vu que les \( f_{\alpha}\) forment une base nous en déduisons \( Q^{-1}AQy=By\). Et vu que \( y\) est un élément quelconque de \( \eK^n\), nous en déduisons l'égalité de matrices
    \begin{equation}        \label{ooWKTYooOJfclT}
        B=Q^{-1}AQ.
    \end{equation}
\end{proof}
Il s'agit bien d'une égalité de matrices, ou à la limite d'applications linéaires sur \( \eK^n\), et non d'une égalité d'application linéaire sur \( V\).

%---------------------------------------------------------------------------------------------------------------------------
\subsection{Changement de base : matrice d'une forme bilinéaire}
%---------------------------------------------------------------------------------------------------------------------------

Soit une forme bilinéaire\footnote{Définition~\ref{DEFooEEQGooNiPjHz}} \( q\colon V\times V\to \eK\) dont la matrice\footnote{Définition~\ref{EQooCUGFooRlKUtu}.} dans la base \( \{ e_i \}\) est \( A\) et celle dans la base \( \{ f_{\alpha} \}\) est \( B\).

Rien n'indique pour l'instant que \( A\) et \( B\) sont les mêmes qu'avant. Au contraire, nous allons voir qu'elles ne sont en général pas les mêmes.

Soit \( x,x'\in V\) de coordonnées \( (x_i)\) et \( (x'_i)\) dans la base \( \{ e_i \}\) et \( (y_{\alpha})\), \( (y'_{\alpha})\) dans la base \( \{ f_{\alpha} \}\). Par définition de la matrice associée à une forme bilinéaire,
\begin{equation}
    q(x,x')=\sum_{ij}A_{ij}x_ix'_j=\sum_{\alpha\beta}B_{\alpha\beta}y_{\alpha}y'_{\beta}.
\end{equation}
En remplaçant les \( x_i\) et \( x'_i\) par leurs valeurs en fonction de \( y_{\alpha}\) et \( y'_{\beta}\),
\begin{subequations}
    \begin{align}
        q(x,x')&=\sum_{ij\alpha\beta}A_{ij}Q_{i\alpha}y_{\alpha}Q_{j\beta}y'_{\beta}\\
        &=\sum_{\alpha\beta}(Q^tAQ)_{\alpha\beta}y_{\alpha}y'_{\beta}
    \end{align}
\end{subequations}
où \( Q^t\) désigne la transposée de la matrice \( Q\) :  \( Q^t_{ij}=Q_{ji}\). Vu que les nombres \( y_{\alpha}\) et \( y'_{\beta}\) sont arbitraires nous déduisons\footnote{Lemme~\ref{LEMooLXAHooPRyHaF}.}
\begin{equation}        \label{EQooZUVTooKjqnJj}
    B=Q^tAQ.
\end{equation}

\begin{remark}
    Notons que cette «loi de transformation» n'est pas la même que celle pour une application linéaire. Ici nous avons \( Q^t\) alors que pour les applications linéaires nous avions \( Q^{-1}\).

    Pour cette raison, tant que nous travaillons avec des bases orthonormées, c'est-à-dire tant que \( Q\) est orthogonale\footnote{Définition~\ref{DefMatriceOrthogonale}.}, nous pouvons confondre une application linéaire avec une application bilinéaire en passant par la matrice. Mais cette identification n'est pas du tout canonique : elle repose sur le fait que les bases soient orthonormées.

    Il en découle que la réduction des endomorphismes et la réduction des formes bilinéaires ne sont pas tout à fait les mêmes théories. Par exemple la pseudo-diagonalisation simultanée (corolaire~\ref{CorNHKnLVA}) est un résultat de réduction de forme bilinéaire et non d'endomorphismes.
    % TODO : lorsque ce sera fait, il faudra approfondir cette remarque.
\end{remark}

%--------------------------------------------------------------------------------------------------------------------------- 
\subsection{Invariance de la trace}
%---------------------------------------------------------------------------------------------------------------------------

\begin{proposition}[\cite{MonCerveau}]      \label{PROPooRMYQooWkEpJJ}
    Soit une application linéaire \( f\). Si la matrice de \( f\) dans une base est \( A\) et est \( B\) dans une autre base, alors
    \begin{equation}
        \trace(A)=\trace(B).
    \end{equation}
\end{proposition}

\begin{proof}
    Les matrices \( A\) et \( B\) sont liées par la proposition \ref{PROPooNZBEooWyCXTw} : \( B=Q^{-1}AQ\) où \( Q\) est la matrice qui lie les vecteurs des deux bases. L'invariance cyclique de la trace donnée en le lemme \ref{LEMooUXDRooWZbMVN} implique que
    \begin{equation}
        \trace(B)=\trace(Q^{-1}AQ)=\trace(QQ^{-1}A)=\trace(A).
    \end{equation}
\end{proof}




%+++++++++++++++++++++++++++++++++++++++++++++++++++++++++++++++++++++++++++++++++++++++++++++++++++++++++++++++++++++++++++
\section{Fonctions}		\label{Sect_fonctions}
%+++++++++++++++++++++++++++++++++++++++++++++++++++++++++++++++++++++++++++++++++++++++++++++++++++++++++++++++++++++++++++

Soient $(V,\| . \|_V)$ et $(W,\| . \|_W)$ deux espaces vectoriels normés, et une fonction $f$ de $V$ dans $W$. Il est maintenant facile de définir les notions de limites et de continuité pour de telles fonctions en copiant les définitions données pour les fonctions de $\eR$ dans $\eR$ en changeant simplement les valeurs absolues par les normes sur $V$ et $W$.

La caractérisation suivante est un recopiage de la définition~\ref{DefOLNtrxB} lorsque la topologie est donnée par des boules.
\begin{proposition}\label{PropHOCWooSzrMjl}
	Soit $f\colon V\to W$ une fonction de domaine \( \Domaine(f)\subset V\) et soit $a$ un point d'accumulation de $\Domaine(f)$.
    La fonction \( f\) admet une limite en $a\in V$ si et seulement s'il existe un élément $\ell\in W$ tel que pour tout $\varepsilon>0$, il existe un $\delta>0$ tel que pour tout $x\in \Domaine(f)$,
    \begin{equation}        \label{EqDefLimzxmasubV}
		0<\| x-a \|_V<\delta\,\Rightarrow\,\| f(x)-\ell \|_W<\varepsilon.
	\end{equation}
	Dans ce cas, nous écrivons $\lim_{x\to a} f(x)=\ell$ et nous disons que $\ell$ est la \defe{limite}{limite} de $f$ lorsque $x$ tend vers $a$.
\end{proposition}

\begin{remark}
    Le fait que nous limitions la formule \eqref{EqDefLimzxmasubV} aux \( x\) dans le domaine de \( f\) n'est pas anodin. Considérons la fonction \( f(x)=\sqrt{x^2-4}\), de domaine \( | x |\geq 2\). Nous avons
    \begin{equation}
        \lim_{x\to 2} \sqrt{x^2-4}=0.
    \end{equation}
    Nous ne pouvons pas dire que cette limite n'existe pas en justifiant que la limite à gauche n'existe pas. Les points \( x<2\) sont hors du domaine de \( f\) et ne comptent dons pas dans l'appréciation de l'existence de la limite.

    Vous verrez plus tard que ceci provient de la \wikipedia{fr}{Topologie_induite}{topologie induite} de \( \eR\) sur l'ensemble \( \mathopen[ 2 , \infty [\).
\end{remark}

%+++++++++++++++++++++++++++++++++++++++++++++++++++++++++++++++++++++++++++++++++++++++++++++++++++++++++++++++++++++++++++
\section{Sous espaces caractéristiques}
%+++++++++++++++++++++++++++++++++++++++++++++++++++++++++++++++++++++++++++++++++++++++++++++++++++++++++++++++++++++++++++

% TODO : lire le blog de Pierre Bernard; en particulier celle-ci : http://allken-bernard.org/pierre/weblog/?p=2299

Lorsqu'un opérateur n'est pas diagonalisable, les valeurs propres jouent quand même un rôle important.

\begin{definition}  \label{DefFBNIooCGbIix}
    Soit \( E\) un \( \eK\)-espace vectoriel  \( f\in\End(E)\). Pour \( \lambda\in \eK\) nous définissons
    \begin{equation}
        F_{\lambda}(f)=\{ v\in E\tq (f-\lambda\mtu)^nv=0, n\in\eN \}
    \end{equation}
    et nous appelons ça un \defe{sous-espace caractéristique}{sous-espace!caractéristique} de \( f\).
\end{definition}
L'espace \( F_{\lambda}(f)\) est l'ensemble de nilpotence de l'opérateur \( f-\lambda\mtu\) et

\begin{lemma}   \label{LemBLPooHMAoyJ}
    L'ensemble \( F_{\lambda}(f)\) est non vide si et seulement si \( \lambda\) est une valeur propre de \( f\). L'espace \( F_{\lambda}(f)\) est invariant sous \( f\).
\end{lemma}

\begin{proof}
    Si \( F_{\lambda}(f)\) est non vide, nous considérons \( v\in F_{\lambda}(f)\) et \( n\) le plus petit entier non nul tel que \( (f-\lambda)^nv=0\). Alors \( (f-\lambda)^{n-1}v\) est un vecteur propre de \( f\) pour la valeur propre \( \lambda\). Inversement si \( v\) est une valeur propre de \( f\) pour la valeur propre \( \lambda\), alors \( v\in F_{\lambda}(f)\).

    En ce qui concerne l'invariance, remarquons que \( f\) commute avec \( f-\lambda\mtu\). Si \( x\in F_{\lambda}(f)\) il existe \( n\) tel que \( (f-\lambda\mtu)^nx=0\). Nous avons aussi
    \begin{equation}
        (f-\lambda\mtu)^nf(x)=f\big( (f-\lambda\mtu)^nx \big)=0,
    \end{equation}
    par conséquent \( f(x)\in F_{\lambda}(f)\).
\end{proof}

\begin{remark}  \label{RemBOGooCLMwyb}
    Toute matrice sur \( \eC\) n'est pas diagonalisable : nous en avons déjà donné une exemple simple en~\ref{ExBRXUooIlUnSx}. Nous en voyons maintenant un moins simple. Considérons en effet l'endomorphisme \( f\) donné par la matrice
    \begin{equation}
        \begin{pmatrix}
            a&    \alpha    &   \beta    \\
            0    &   a    &   \gamma    \\
            0    &   0    &   b
        \end{pmatrix}
    \end{equation}
    où \( a\neq b\), \( \alpha\neq 0\), \( \beta\) et \( \gamma\) sont des nombres complexes quelconques.
    Son polynôme caractéristique est
    \begin{equation}
        \chi_f(\lambda)=(a-\lambda)^2(b-\lambda),
    \end{equation}
    et les valeurs propres sont donc \( a\) et \( b\). Nous trouvons les vecteurs propres pour la valeur \( a\) en résolvant
    \begin{equation}
        \begin{pmatrix}
            a    &   \alpha    &   \beta    \\
            0    &   a    &   \gamma    \\
            0    &   0    &   b
        \end{pmatrix}\begin{pmatrix}
            x    \\
            y    \\
            z
        \end{pmatrix}=\begin{pmatrix}
            ax    \\
            ay    \\
            az
        \end{pmatrix}.
    \end{equation}
    L'espace propre \( E_a(f)\) est réduit à une seule dimension générée par \( (1,0,0)\). De la même façon l'espace propre correspondant à la valeur propre \( b\) est donné par
    \begin{equation}
        \begin{pmatrix}
            \frac{1}{ b-a }\left( \beta+\frac{ \alpha\gamma }{ b-a } \right)    \\
            \frac{ \gamma }{ b-a }    \\
            1
        \end{pmatrix}.
    \end{equation}
    Il n'y a donc pas trois vecteurs propres linéairement indépendants, et l'opérateur \( f\) n'est pas diagonalisable.

    Par contre nous pouvons voir que
    \begin{equation}
        (f-\alpha\mtu)^2\begin{pmatrix}
             0   \\
            1    \\
            0
        \end{pmatrix}=
        \begin{pmatrix}
            a    &   \alpha    &   \beta    \\
            0    &   a    &   \gamma    \\
            0    &   0    &   b
        \end{pmatrix}
        \begin{pmatrix}
            \alpha    \\
            0    \\
            0
        \end{pmatrix}-\begin{pmatrix}
            a\alpha    \\
            0    \\
            0
        \end{pmatrix}=\begin{pmatrix}
            0    \\
            0    \\
            0
        \end{pmatrix},
    \end{equation}
    de telle sorte que le vecteur \( (0,1,0)\) soit également dans l'espace caractéristique \( F_a(f)\).

    Dans cet exemple, la multiplicité algébrique de la racine \( a\) du polynôme caractéristique vaut \( 2\) tandis que sa multiplicité géométrique vaut seulement \( 1\).
\end{remark}

%---------------------------------------------------------------------------------------------------------------------------
\subsection{Théorèmes de décomposition}
%---------------------------------------------------------------------------------------------------------------------------

%TODO : Je crois qu'on peut remplacer l'hypothèse de corps algébriquement clos par le polynôme caractéristique scindé.
\begin{theorem}[Théorème spectral, décomposition primaire]\index{théorème!spectral}     \label{ThoSpectraluRMLok}
    Soit \( E\) espace vectoriel de dimension finie sur le corps algébriquement clos \( \eK\) et \( f\in\End(E)\). Alors
    \begin{equation}    \label{EqCTFHooBSGhYK}
        E=F_{\lambda_1}(f)\oplus\ldots\oplus F_{\lambda_k}(f)
    \end{equation}
    où la somme est sur les valeurs propres distinctes de \( f\).

    Les projecteurs sur les espaces caractéristique forment un système complet et orthogonal.
\end{theorem}
\index{décomposition!primaire}
\index{décomposition!spectrale}
\index{décomposition!sous-espaces caractéristiques}

\begin{proof}
    Soit \( P\) le polynôme caractéristique de \( f\) et une décomposition
    \begin{equation}
        P=(f-\lambda_1)^{\alpha_1}\ldots(f-\lambda_r)^{\alpha_r}
    \end{equation}
    en facteurs irréductibles. La le théorème de noyaux (\ref{ThoDecompNoyayzzMWod}) nous avons
    \begin{equation}        \label{EqDeFVSaYv}
        E=\ker(f-\lambda_1)^{\alpha_1}\oplus\ldots\oplus\ker(f-\lambda_r)^{\alpha_r}.
    \end{equation}
    Les projecteurs sont des polynômes en \( f\) et forment un système orthogonal. Il nous reste à prouver que \( \ker(f-\lambda_i)^{\alpha_i}=F_{\lambda_i}(f)\). L'inclusion
    \begin{equation}    \label{EqzmNxPi}
        \ker(f-\lambda_i)^{\alpha_i}\subset F_{\lambda_i}(f)
    \end{equation}
    est évidente. Nous devons montrer l'inclusion inverse. Prouvons que la somme des \( F_{\lambda_i}(f)\) est directe. Si \( v\in F_{\lambda_i}(f)\cap F_{\lambda_j}(f)\), alors il existe \( v_1=(f-\lambda_i)^nv\neq 0\) avec \( (f-\lambda_i)v_1=0\). Étant donné que \( (f-\lambda_i)\) commute avec \( (f-\lambda_j)\), ce \( v_1\) est encore dans \( F_{\lambda_j}(f)\) et par conséquent il existe \( w=(f-\lambda_j)^mv_1\) non nul tel que
    \begin{subequations}
        \begin{numcases}{}
            (f-\lambda_i)w=0\\
            (f-\lambda_j)w=0.
        \end{numcases}
    \end{subequations}
    Ce \( w\) serait donc un vecteur propre simultané pour les valeurs propres \( \lambda_i\) et \( \lambda_j\), ce qui est impossible parce que les espaces propres sont linéairement indépendants. Les espaces \( F_{\lambda_i}\) sont donc en somme directe et
    \begin{equation}
        \sum_i\dim F_{\lambda_i}(f)\leq \dim E.
    \end{equation}
    En tenant compte de l'inclusion \eqref{EqzmNxPi} nous avons même
    \begin{equation}
        \dim E=\sum_i\dim\ker(f-\lambda_i)^{\alpha_i}\leq\sum_i F_{\lambda_i}(f)\leq \dim E.
    \end{equation}
    Par conséquent nous avons \( \dim\ker(f-\lambda_i)^{\alpha_i}=\dim F_{\lambda_i}(f)\) et l'égalité des deux espaces.
\end{proof}


\begin{probleme}
    Dans le cas où le corps n'est pas algébriquement clos, il paraît qu'il faut remplacer «diagonalisable» par «semi-simple».
\end{probleme}
%TODO : peut-être qu'il y a la réponse dans http://www.math.jussieu.fr/~romagny/agreg/dvt/endom_semi_simples.pdf

Si l'espace vectoriel est sur un corps algébriquement clos, alors les endomorphismes semi-simples\footnote{Définition~\ref{DEFooBOHVooSOopJN}.} sont les endomorphismes diagonaux.


%TODO : Je crois qu'on peut remplacer l'hypothèse de corps algébriquement clos par le polynôme caractéristique scindé.
\begin{theorem}[Décomposition de Dunford] \label{ThoRURcpW}
    Soit \( E\) un espace vectoriel sur le corps algébriquement clos \( \eK\) et \( u\in\End(E)\) un endomorphisme de \( E\).

    \begin{enumerate}
        \item

            L'endomorphisme \( u\) se décompose de façon unique sous la forme
            \begin{equation}
                u=s+n
            \end{equation}
            où \( s\) est diagonalisable, \( n\) est nilpotent et \( [s,n]=0\).
        \item
            Les endomorphismes \( s\) et \( n\) sont des polynômes en \( u\) et commutent avec \( u\).
        \item   \label{ItemThoRURcpWiii}
            Les parties \( s\) et \( n\) sont données par
            \begin{subequations}
                \begin{align}
                    s&=\sum_i\lambda_ip_i\\
                    n&=\sum_i(s-\lambda_i\mtu)p_i
                \end{align}
            \end{subequations}
            où les sommes sont sur les valeurs propres distinctes\footnote{C'est-à-dire sur les sous-espaces caractéristiques.} de \( f\) et où \( p_i\colon E\to F_{\lambda_i}(u)\) est la projection de \( E\) sur \( F_{\lambda_i}(u)\).
    \end{enumerate}
\end{theorem}
\index{décomposition!Dunford}
\index{Dunford!décomposition}
\index{réduction!d'endomorphisme}
\index{endomorphisme!sous-espace stable}
\index{polynôme!d'endomorphisme!décomposition de Dunford}
\index{endomorphisme!diagonalisable!Dunford}
\index{endomorphisme!nilpotent!Dunford}
%TODO : comprendre comment on calcule des exponentielles de matrices avec Dunford.

\begin{proof}
    Le théorème spectral~\ref{ThoSpectraluRMLok} nous indique que
    \begin{equation}
        E=\bigoplus_iF_{\lambda_i}(f).
    \end{equation}
    Nous considérons l'endomorphisme \( s\) de \( E\) qui consiste à dilater d'un facteur \( \lambda\) l'espace caractéristique \( F_{\lambda}(f)\) :
    \begin{equation}
        s=\sum_i\lambda_ip_i
    \end{equation}
    où \( p_i\colon E\to F_{\lambda_i}(u)\) est la projection de \( E\) sur \( F_{\lambda_i}(u)\).

    Nous allons prouver que \( [s,f]=0\) et \( n=f-s\) est nilpotent. Cela impliquera que \( [s,n]=0\).

    Si \( x\in F_{\lambda}(f)\), alors nous avons \( sf(x)=\lambda f(x)\) parce que \( f(x)\in F_{\lambda}(f)\) tandis que \( fs(x)=f(\lambda x)=\lambda f(x)\). Par conséquent \( f\) commute avec \( s\).

    Pour montrer que \( f-s\) est nilpotent, nous en considérons la restriction
    \begin{equation}
        f-s\colon F_{\lambda}(f)\to F_{\lambda}(f).
    \end{equation}
    Cet opérateur est égal à \( f-\lambda\mtu\) et est par conséquent nilpotent.

    Prouvons à présent l'unicité. Soit \( u=s'+n'\) une autre décomposition qui satisfait aux conditions : \( s'\) est diagonalisable, \( n'\) est nilpotent et \( [n',s']=0\). Commençons par prouver que \( s'\) et \( n'\) commutent avec \( u\). En multipliant \( u=s'+n'\) par \( s'\) nous avons
    \begin{equation}
        s'u=s'^2+s'n'=s'^2+n's'=(s'+n')s'=us',
    \end{equation}
    par conséquent \( [u,s']=0\). Nous faisons la même chose avec \( n'\) pour trouver \( [u,n']=0\). Notons que pour obtenir ce résultat nous avons utilisé le fait que \( n'\) et \( s'\) commutent, mais pas leur propriétés de nilpotence et de diagonalisibilité.


    Si \( s'+n'=s+n\) est une autre décomposition, \( s'\) et \( n'\) commutent avec \( u\), et par conséquent avec tous les polynômes en \( u\). Ils commutent en particulier avec \( n\) et \( s\). Les endomorphismes \( s\) et \( s'\) sont alors deux endomorphismes diagonalisables qui commutent. Par la proposition~\ref{PropGqhAMei}, ils sont simultanément diagonalisables. Dans la base de simultanée diagonalisation, la matrice de l'opérateur \( s'-s=n-n'\) est donc diagonale. Mais \( n-n'\) est également nilpotent, en effet si \( A\) et \( B\) sont deux opérateurs nilpotents,
    \begin{equation}
        (A+B)^n=\sum_{k=0}^n\binom{k}{n}A^kB^{n-k}.
    \end{equation}
    Si \( n\) est assez grand, au moins un parmi \( A^k\) ou \( B^{n-k}\) est nul.

    Maintenant que \( n-n'\) est diagonal et nilpotent, il est nul et \( n=n'\). Nous avons alors immédiatement aussi \( s=s'\).

\end{proof}

%---------------------------------------------------------------------------------------------------------------------------
\subsection{Diverses conséquences}
%---------------------------------------------------------------------------------------------------------------------------

\begin{theorem}
    Soit une matrice \( A\in \eM(n,\eC)\). On a que la suite \( (A^kx)\) tend vers zéro pour tout \( x\) si et seulement si \( \rho(A)<1\) où \( \rho(A)\)\index{rayon!spectral} est le rayon spectral de $A$
\end{theorem}
\index{décomposition!Dunford!exponentielle de matrice}

\begin{proof}
    Dans le sens direct, il suffit de prendre comme \( x\), un vecteur propre de \( A\). Dans ce cas nous avons \( A^kx=\lambda^kx\). Mais \( \lambda^kx\) ne tend vers zéro que si \( \lambda<1\). Donc toutes les valeurs propres de \( A\) doivent être plus petite que \( 1\) et \( \rho(A)<1\).

    Pour l'autre sens nous utilisons la décomposition de Dunford (théorème~\ref{ThoRURcpW}) : il existe une matrice inversible \( P\) telle que
    \begin{equation}
        A=P^{-1}(D+N)P
    \end{equation}
    où \( D\) est diagonale, \( N\) est nilpotente et \( [D,N]=0\). Étant donné que \( D+N\) est triangulaire, son polynôme caractéristique que
    \begin{equation}
        \chi_{D+N}(\lambda)=\prod_i D_{ii}-\lambda.
    \end{equation}
    Par similitude, c'est le même polynôme caractéristique que celui de \( A\) et nous savons alors que la diagonale de \( D\) contient les valeurs propres de \( A\).

    Par ailleurs nous avons
    \begin{subequations}
        \begin{align}
            A^k&=P^{-1}(D+N)^kP\\
            &=P^{-1}\sum_{j=0}^k{j\choose k}D^{j-k}N^jP\\
            &=P^{-1}\sum_{j=0}^{n-1}{j\choose k}D^{j-k}N^jP
        \end{align}
    \end{subequations}
    où nous avons utilité le fait que \( D\) et \( N\) commutent ainsi que \( N^{n-1}=0\) parce que \( N\) est nilpotente. Nous utilisons la norme matricielle usuelle, pour laquelle \( \| D \|=\rho(D)=\rho(A)\). Nous avons alors
    \begin{equation}
        \| (D+N)^k \|\leq \sum_{j=0}^k{j\choose k}\rho(D)^{k-j}\| N \|^j.
    \end{equation}
    Du coup si \( \rho(D)<1\) alors \( \| (D+N)^k \|\to 0\) (et c'est même un si et seulement si).
\end{proof}

Une application de la décomposition de Jordan est l'existence d'un logarithme pour les matrices. La proposition suivant va d'une certaine manière donner un logarithme pour les matrices inversibles complexes. Dans le cas des matrices réelles \( m\) telles que \( \| m-\mtu \|<1\), nous donnerons au lemme~\ref{LemQZIQxaB} une formule pour le logarithme sous forme d'une série; ce logarithme sera réel.

%---------------------------------------------------------------------------------------------------------------------------
\subsection{Valeurs singulières}
%---------------------------------------------------------------------------------------------------------------------------

\begin{definition}
    Soit \( M\) une matrice \( m\times n\) sur \( \eK\) (\( \eK\) est \( \eR\) ou \( \eC\)). Un nombre réel \( \sigma\) est une \defe{valeur singulière}{valeur!singulière} de \( M\) s'il existent des vecteurs unitaires \( u\in \eK^m\), \( v\in \eK^n\) tels que
    \begin{subequations}
        \begin{align}
            Mv&=\sigma u\\
            M^*u&=\sigma v.
        \end{align}
    \end{subequations}
\end{definition}

\begin{theorem}[Décomposition en valeurs singulières]
    Soit \( M\in \eM(m\times n,\eK)\) où \( \eK=\eR,\eC\). Alors \( M\) se décompose en
    \begin{equation}
        M=ADB
    \end{equation}
    où
    il existe deux matrices unitaires \( A\in \gU(m\times m)\), \( B\in \gU(n\times n)\) et une matrice (pseudo)diagonale \( D\in \eM(m\times n)\) tels que
    \begin{enumerate}
        \item
            \( A\in\gU(m\times m)\), \( B\in\gU(n\times n)\) sont deux matrices unitaires;,
        \item
            \( D\) est (pseudo)diagonale,
        \item
            les éléments diagonaux de \( D\) sont les valeurs singulières de \( M\),
        \item
            le nombre d'éléments non nuls sur la diagonale de \( D\) est le rang\footnote{Définition~\ref{DefALUAooSPcmyK}.} de \( M\).
    \end{enumerate}
\end{theorem}

\begin{corollary}
    Soit \( M\in \eM(n,\eC)\). Il existe un isomorphisme \( f\colon \eC^n\to \eC^n\) tel que \( fM\) soit autoadjoint.
\end{corollary}

\begin{proof}
    Si \( M=ADB\) est la décomposition de \( M\) en valeurs singulières, alors nous pouvons prendre \( f=\overline{ B }^tA^{-1}\) qui est une matrice inversible. Pour la vérification que ce \( f\) répond bien à la question, ne pas oublier que \( D\) est réelle, même si \( M\) ne l'est pas.
\end{proof}
